\begin{proposition}\label{prop:decic-no117}
An elliptic curve over a number field of degree $10$ has no point
of order $117$.
\end{proposition}
\begin{proof}
Suppose that $E/K$ has a point $Q$ of order $117$, with
$[K:\mathbb Q]=10$. Let $D$ be the rational trace divisor of
$(E,3Q)$ on $X=X_1(39)$, including multiplicities if its moduli
field is smaller than $K$.
We use the degree-independent $X_1(39)$ certificate. The curve has
genus $33$, rational gonality $14$, and
\[
 J_X(\mathbb Q)=[7,31122,3236688],\qquad
 J_X(\mathbb Q)[2^\infty]=[2,16].
\]
The Jacobian statement follows from
\cite[Theorem~3.1(2), Theorem~4.13 and Table~2]{DEHMZ};
the gonality is given in \cite[Table~1]{DvH}.
Specialization at $2$ is injective on $J_X(\mathbb Q)$.
We recall its independent certification: the order-$3$ diamond
quotient has genus $9$ and rational cuspidal $2$-primary reduction
image of order $32$. Pullback bounds its rational $2$-primary
torsion by the source group $[2,16]$, also of order $32$.
The order-$3$ diamond acts trivially on this source group, since
its automorphism group has order $32$.
On the source rational $2$-primary group the norm identity
$\pi^*\pi_*=1+h+h^2=[3]$, where $h$ is the order-$3$ diamond,
then proves injectivity; odd torsion is automatically
injective. The quotient model, actual cuspidal relation lattice,
and the order-$32$ reduction image are all checked in the retained
certificate.

The global and finite cusp signatures both equal
\[
                         1^{12}2^6 12^2.
\]
Hence every cuspidal divisor of degree at most $10$ in the special
fibre lifts to a rational cuspidal divisor of the same degree.
If $\overline D$ is entirely cuspidal, let $C$ be this lift.
Injectivity gives $D\sim C$, and gonality $14>10$ gives $D=C$,
contrary to the noncuspidal support of $D$.

At every noncuspidal reduction of $(E,3Q)$ the fine level-$39$
model gives good reduction of $E$. The prime-to-$2$ point $Q$
retains order $117$. The independent complete census over
$\mathbb F_{2^f}$, $1\leq f\leq10$, shows that the only possible
residue degrees are $8$ and $9$, with group orders $234$ and
$468$, respectively. In particular no new degree-$10$ reduction
occurs. This last assertion can also be seen from the Hasse
interval $[961,1089]$: its only multiple of $117$ is $1053$,
whose trace $-28$ is neither ordinary nor supersingular in
characteristic $2$.
If the actual residue field of $E$ has degree $f=8$ or $9$,
the degree $g$ of its marked $39$-level place divides $f$.
Every proper divisor is at most $4$, where the Hasse bound
excludes order $39$. Thus $g=f$.
It follows that $\overline D=P+C$, where $P$ is one noncuspidal
place of degree $8$ or $9$, and $C$ is a cuspidal divisor of degree $2$ or
$1$. Both complements lift globally.

Put $A_5=T_5-1-5\langle5\rangle$, using the point-level Hecke
correspondence. Eichler--Shimura and torsion injectivity at $5$
give $A_5J_X(\mathbb Q)=0$. The cusp boundary formula
\[
 T_5=\operatorname{Fr}_5+
                5\langle5\rangle\operatorname{Fr}_5^{-1}
\]
shows that $A_5$ annihilates every rational cuspidal divisor
literally. Hence $A_5P$ must be principal.
There are five octic diamond orbits. The retained exact
nonprincipality tests exclude four of them; the fifth is the
reduction of the known closed octic divisor $D_8$.
In this case injectivity and gonality give
$D=hD_8+C$ for a rational diamond $h$ and the rational lift of
$C$, contradicting noncuspidal support.

It remains to consider degree $9$.
The complete census has exactly three curves of order $468$,
whose $j$-invariants form a degree-$3$ orbit.
Their rational $39$-torsion is cyclic since $\gcd(39,511)=1$.
Thus they give exactly $3\varphi(39)/2=36$ geometric marked
points. The known closed nonic divisor $D_9$ has four distinct
closed diamond translates, giving all $36$ points.
Independent exact finite-field arithmetic binds those four
reductions to the complete census.
Consequently $\overline D=\overline{hD_9}+C$ for some rational
diamond $h$. Injectivity and gonality again give $D=hD_9+C$,
contradicting noncuspidal support.
\end{proof}
